\section{Day 5 -- Tuning the proposal variance $\Sigma$}

The Metropolis--Hastings sampler used in Part~3 relies on a symmetric Gaussian
random-walk proposal
\begin{equation}
    h(x,x') \;=\; \mathcal{N}\!\bigl(x'\,\big|\,x,\,\Sigma^{2}\bigr),
\end{equation}
applied coordinate-wise to one of $(\theta,\psi)$ per step. The step size
$\Sigma$ does not affect the stationary distribution -- as long as the chain
is irreducible the target is still the posterior -- but it does control how
efficiently the chain explores the state space. We illustrate this by running
the sampler with $\Sigma\in\{0.01,\;1,\;50\}$ on the same two data sets used
in Part~3: $n=10$ observations (where we use the posterior on the original
scale) and $n=1000$ observations (where we use the log-posterior for
numerical stability). The two extremes bracket the well-tuned case so the
trade-off is easy to see.

\subsection*{Too small a step ($\Sigma=0.01$)}

With a very small proposal variance almost every move is accepted, but each
move barely changes the state. The chain wanders in tiny increments and
takes far too long to traverse the support of the posterior. The result is
visible in Figure~\ref{fig:sigma_small}: the empirical density is patchy and
multimodal-looking, with isolated lumps that the kernel density estimator
turns into spurious contour rings. These ``modes'' are not features of the
posterior -- they are samples from a single Markov chain that has not had
time to mix between regions of the same well.

\begin{figure}[H]
    \centering
    \begin{subfigure}[t]{0.48\linewidth}
        \includegraphics[width=\linewidth]{plots/day5/posterior_samples_xobs10sigma_001.png}
        \caption{$n=10$ observations.}
    \end{subfigure}
    \hfill
    \begin{subfigure}[t]{0.48\linewidth}
        \includegraphics[width=\linewidth]{plots/day5/log_posterior_samples_sigma_001.png}
        \caption{$n=1000$ observations (log-posterior).}
    \end{subfigure}
    \caption{Posterior samples with $\Sigma=0.01$. The chain takes tiny
    steps; high acceptance but very slow mixing. The empirical density is
    rough and shows spurious local structure.}
    \label{fig:sigma_small}
\end{figure}

\subsection*{Well-tuned step ($\Sigma=1$)}

A proposal standard deviation on the same order of magnitude as the
posterior standard deviation is close to optimal. The acceptance rate sits
comfortably between the two extremes, the chain visits the bulk of the
posterior in a reasonable number of steps, and the empirical density is
smooth and unimodal in both data regimes
(Figure~\ref{fig:sigma_one}). The $n=1000$ panel is tightly concentrated
around the true $(\theta,\psi)$, as expected once the likelihood dominates
the prior.

\begin{figure}[H]
    \centering
    \begin{subfigure}[t]{0.48\linewidth}
        \includegraphics[width=\linewidth]{plots/day5/posterior_samples_xobs10.png}
        \caption{$n=10$ observations.}
    \end{subfigure}
    \hfill
    \begin{subfigure}[t]{0.48\linewidth}
        \includegraphics[width=\linewidth]{plots/day5/log_posterior_samples.png}
        \caption{$n=1000$ observations (log-posterior).}
    \end{subfigure}
    \caption{Posterior samples with $\Sigma=1$. Smooth, unimodal empirical
    densities; the chain mixes well across the support.}
    \label{fig:sigma_one}
\end{figure}

\subsection*{Too large a step ($\Sigma=50$)}

When $\Sigma$ is much larger than the posterior standard deviation, almost
every proposal lands far in the tail of the target where the density is
negligible compared to the current state, and the acceptance ratio
collapses. The chain stays put for long stretches, occasionally jumping to a
random distant point that happens to be accepted. The samples in
Figure~\ref{fig:sigma_large} are sparse and ragged: the histograms are
spiky, the KDE contours fragment into noise, and the $n=1000$ panel in
particular shows so few distinct accepted states that the kernel density
estimate is essentially uninformative.

\begin{figure}[H]
    \centering
    \begin{subfigure}[t]{0.48\linewidth}
        \includegraphics[width=\linewidth]{plots/day5/posterior_samples_xobs10_sigma50.png}
        \caption{$n=10$ observations.}
    \end{subfigure}
    \hfill
    \begin{subfigure}[t]{0.48\linewidth}
        \includegraphics[width=\linewidth]{plots/day5/log_posterior_samples_sigma50.png}
        \caption{$n=1000$ observations (log-posterior).}
    \end{subfigure}
    \caption{Posterior samples with $\Sigma=50$. Almost every proposal is
    rejected; the chain spends long stretches at the same state and the
    resulting density is sparse and noisy.}
    \label{fig:sigma_large}
\end{figure}

\subsection*{Discussion}

The three regimes illustrate the classical Metropolis--Hastings tuning
trade-off:
\begin{itemize}
    \item $\Sigma$ too small $\Rightarrow$ high acceptance, low movement,
        the chain is locally Brownian and fails to traverse the posterior in
        the available budget.
    \item $\Sigma$ too large $\Rightarrow$ low acceptance, the chain is
        effectively frozen and only occasionally jumps to a usable state.
    \item $\Sigma$ comparable to the posterior standard deviation
        $\Rightarrow$ acceptance rate roughly in the $0.2$--$0.5$ range, and
        the chain explores efficiently.
\end{itemize}
The effect is more pronounced for $n=1000$ than for $n=10$ because the
posterior at $n=1000$ is much narrower: a step of $\Sigma=50$ overshoots the
posterior by orders of magnitude, while a step of $\Sigma=0.01$ now needs
thousands of accepted moves to cross even a single posterior standard
deviation. The well-tuned $\Sigma$ therefore implicitly depends on $n$ --
in a practical implementation one would either tune $\Sigma$ adaptively
during a burn-in phase or scale it with an estimate of the posterior
covariance.
